%! AP Statistics — Unit 8, Lesson 8.2 — Warm-Up (Answer Key)
\documentclass[10pt]{article}
\usepackage{apstats-article}
\usepackage{apstats-key}

\begin{document}

\pageheader{Unit 8, Lesson 8.2}{Warm-Up --- Answer Key}
\namedateperiod

\vspace{0.15in}

A school counselor claims that students choose among three lunch options in equal
proportions ($1/3$ each). In a random sample of 90 students, 38 chose pizza, 29 chose
salad, and 23 chose pasta.

\begin{enumerate}

  \item State H$_0$ and H$_a$ \textbf{in words} for a test of whether the three lunch
        options are chosen in equal proportions.

        H$_0$: \ans{Pizza, salad, and pasta are each chosen by $1/3$ of students
        ($p_\text{pizza} = p_\text{salad} = p_\text{pasta} = 1/3$).}

        \vspace{0.05in}

        H$_a$: \ans{At least one lunch option is chosen by a proportion different from $1/3$.}

        \vspace{0.05in}

  \item If H$_0$ is true (each proportion $= 1/3$), how many students would you
        \textbf{expect} to choose each option in a sample of 90?

        \vspace{0.05in}
        $E = n \cdot p_0 = 90 \cdot \dfrac{1}{3} = $ \ans{30 students for each option.}

        \vspace{0.1in}

  \item From Unit 6, what conditions must be met to conduct valid inference from a
        random sample? List at least two.

        \begin{enumerate}[label=\alph*.]
          \item \ans{Random sample (ensures independence of observations).}
          \item \ans{Large enough counts / large sample condition (ensures the sampling
                distribution has an appropriate shape).}
        \end{enumerate}

\end{enumerate}

\end{document}
